% Exam copy: my scribes (../scribes/tex) with practice questions after each week.
% Generated from scribes/main.tex; build from this folder, three passes.
% The lectures are grouped by week (\newweek{number}{dates}); each week has a Tuesday and a
% Thursday lecture. Every lecture day has its own file in tex/ and starts with
% \newlecture{weekday, date}{topic}.
% scribe.sty is an unmodified copy of the course file.

\documentclass[a4paper,11pt,final]{article}
\makeatletter
\def\input@path{{../scribes/}}
\makeatother
\usepackage{scribecompat} % must come before scribe
\usepackage{scribe}
\usepackage{scribeextras} % \newlecture, dashed exercise boxes, \solutionspace, inline TikZ

\theoremstyle{plain}
\newtheorem{claim}[theorem]{Claim}
\crefname{claim}{Claim}{Claims}
\Crefname{claim}{Claim}{Claims}


% Practice questions after each week (not part of the submitted scribes).
\theoremstyle{definition}
\newcounter{practiceweek}
\newtheorem{pproblem}{Practice Problem}
\renewcommand{\thepproblem}{\arabic{practiceweek}.\arabic{pproblem}}
\crefname{pproblem}{Practice Problem}{Practice Problems}
\Crefname{pproblem}{Practice Problem}{Practice Problems}
\newcommand{\scr}[1]{\cref{#1}, p.~\pageref{#1}}
\newcommand{\uses}[1]{\noindent\textit{Scribes:} #1.\par}
\AtBeginEnvironment{wsolution}{\setlength{\parskip}{0.35em}}
\makeatletter
\newcommand{\practiceweek}[1]{%
    \par\FloatBarrier
    \vspace{0.4in}%
    \needspace{2.5in}%
    \setcounter{practiceweek}{#1}%
    \setcounter{pproblem}{0}%
    \phantomsection
    \addcontentsline{toc}{lecture}{Practice questions, week #1}%
    \hrule
    \vspace{0.1in}
    \noindent\textbf{Practice questions for week #1} (likely exam questions, with solutions)\par
    \vspace{0.1in}
    \hrule
    \vspace{0.1in}
}
\makeatother

\begin{document}

%%%%%%%%%%%%%%%%%%%%%%%%%%%%%%%%%%%%%%%%%%%%%%%%%%%%%%%%%%%%%%%%%%%%%%%%%%%
% Document begins here %%%%%%%%%%%%%%%%%%%%%%%%%%%%%%%%%%%%%%%%%%%%%%%%%%%%
%%%%%%%%%%%%%%%%%%%%%%%%%%%%%%%%%%%%%%%%%%%%%%%%%%%%%%%%%%%%%%%%%%%%%%%%%%%

\begin{center}
\textbf{\large Quantum Computing} \\

Professor Gopal Pandurangan\\
\vspace{0.1in}
\end{center}

\hrule

\vspace{0.1in}

\textbf{Scribe Author:} Armaan Amatya

\vspace{0.1in}

\hrule

\vspace{0.1in}

\tableofcontents

%%%%%%%%%%%%%%%%%%%%%%%%%%%%%%%%%%%%%%%%%%%%%%%%%%%%%%%%%%%%%%%%%%%%%%%%%%%
% Lectures by week, in order. Append new weeks and lectures at the end. %%%
%%%%%%%%%%%%%%%%%%%%%%%%%%%%%%%%%%%%%%%%%%%%%%%%%%%%%%%%%%%%%%%%%%%%%%%%%%%

\newweek{1}{August 24--28, 2026}

\newlecture{Tuesday, August 25, 2026}{Efficient Algorithms, Random Sampling, Lower Bound for
    Deterministic Search, Probability Basics}
\input{tex/lecture_2026-08-25}

\newlecture{Thursday, August 27, 2026}{Search and Satisfiability, Expectation, Bernoulli, Binomial and
    Geometric Random Variables, Conditional Expectation}
\input{tex/lecture_2026-08-27}

\practiceweek{1}

\begin{exercisebox}
\begin{pproblem}[Membership needs $n$ queries]\label{pp:membership}
    The scribe shows that finding $x$ needs $n - 1$ queries when $x$ is promised to be in the array.
    Show that deciding whether $x$ is in the array at all needs $n$ queries.
\end{pproblem}
\tcblower
\begin{wsolution}
    \uses{\scr{thm:detsearch-0825}}
    \leavevmode
    \begin{enumerate}
        \item Fix a deterministic algorithm. The adversary answers every query with ``not $x$''.
        \item Suppose the algorithm stops after at most $n - 1$ queries. Then some slot $j$ was never
              queried.
        \item Two arrays agree with all the answers: one has $x$ in slot $j$, the other has no $x$
              anywhere.
        \item The algorithm saw the same answers on both, so it gives the same output on both. It is
              wrong on one.
        \item So a correct algorithm makes $n$ queries.
    \end{enumerate}
    The promise is what saved one query on Aug 25: with a promise, $n - 1$ answers of ``no'' leave
    only one possible slot.
\end{wsolution}
\end{exercisebox}

\begin{exercisebox}
\begin{pproblem}[Search with $k$ copies of $x$]\label{wp:kcopies}
    An array of $n$ slots contains $x$ in exactly $k$ slots, $1 \leq k < n$.
    \begin{enumerate}[label=(\alph*)]
        \item Random sampling queries uniformly random slots (with replacement) until it finds $x$.
              What is the expected number of queries?
        \item Show that a deterministic algorithm needs $n - k$ queries in the worst case.
    \end{enumerate}
\end{pproblem}
\tcblower
\begin{wsolution}
    \uses{\scr{thm:geometric-0827}; \scr{thm:detsearch-0825}}
    (a) Each query finds $x$ with probability $p = \sfrac{k}{n}$, independently of the others, so the
    number of queries is geometric with $p = \sfrac{k}{n}$. Its expectation is
    $\sfrac{1}{p} = \sfrac{n}{k}$.

    (b) The adversary answers ``not $x$'' to every query. Suppose the algorithm stops after
    $q < n - k$ queries and outputs a slot $j$. The array has $n - k$ slots without $x$, and only
    $q$ of them have been used up by the answers, so at least one more slot without $x$ remains. If
    $j$ was queried, the answer at $j$ was ``not $x$'' and the output is wrong. Otherwise the
    adversary puts the remaining empty slot at $j$ and the $k$ copies of $x$ in other unqueried
    slots; this agrees with all the answers, and the output is wrong. After $n - k$ answers of
    ``no'', the $k$ unqueried slots all contain $x$, so $n - k$ queries suffice.

    For $k = 1$ this is the $n - 1$ of the scribe.
\end{wsolution}
\end{exercisebox}


\newweek{2}{August 31 -- September 4, 2026}

\newlecture{Tuesday, September 1, 2026}{Query Complexity of Search, Coupon Collector Problem,
    Lower Bound for Randomized Algorithms, Chernoff Bounds}
\input{tex/lecture_2026-09-01}

\newlecture{Thursday, September 3, 2026}{Markov, Chebyshev and Chernoff Inequalities, Search,
    Birthday Paradox}
\input{tex/lecture_2026-09-03}

\practiceweek{2}

\begin{exercisebox}
\begin{pproblem}[$n$ balls into $n$ bins]\label{pp:balls}
    Throw $n$ balls independently and uniformly into $n$ bins, $n \geq 4$. Show that with probability
    at least $1 - \sfrac{1}{n^2}$ no bin receives $3 \log_2 n$ balls or more.
\end{pproblem}
\tcblower
\begin{wsolution}
    \uses{\scr{thm:chernoff-0901}}
    \leavevmode
    \begin{enumerate}
        \item Fix a bin and let $X$ be its number of balls. $X$ is a sum of $n$ independent
              indicators, each $1$ with probability $\sfrac{1}{n}$, so $\mu = 1$.
        \item Take $R = 3 \log_2 n$. For $n \geq 4$ we have $R \geq 6 = 6\mu$, so the third Chernoff
              form applies.
        \item $\prob{X \geq 3 \log_2 n} \leq 2^{-3 \log_2 n} = \sfrac{1}{n^3}$.
        \item Union bound over the $n$ bins:
              $\prob{\text{some bin has at least } 3 \log_2 n \text{ balls}} \leq n \cdot \frac{1}{n^3} = \frac{1}{n^2}$. \qedhere
    \end{enumerate}
\end{wsolution}
\end{exercisebox}

\begin{exercisebox}
\begin{pproblem}[Coupon collector with a different number of trials]\label{pp:coupon}
    There are $n$ coupons and each trial gives a uniformly random one. Bound the probability that
    some coupon is missing after
    \begin{enumerate}[label=(\alph*)]
        \item $2n \ln n$ trials,
        \item $n \ln n + cn$ trials.
        \item What happens after only $n$ trials?
    \end{enumerate}
\end{pproblem}
\tcblower
\begin{wsolution}
    \uses{\scr{thm:whp}; \scr{thm:linearity-0827}}
    A fixed coupon is missed by one trial with probability $1 - \sfrac{1}{n}$, so it is missing
    after $m$ trials with probability $\left(1 - \frac{1}{n}\right)^m \leq e^{-m/n}$.

    (a) $e^{-2 \ln n} = \sfrac{1}{n^2}$ per coupon, so at most $n \cdot \frac{1}{n^2} = \frac{1}{n}$
    by the union bound.

    (b) $e^{-\ln n - c} = \sfrac{e^{-c}}{n}$ per coupon, so at most $e^{-c}$ in total. For $c = 3$
    this is about $0.05$.

    (c) The expected number of missing coupons is
    $n\left(1 - \frac{1}{n}\right)^n \approx \sfrac{n}{e}$, by linearity of expectation. About
    $37\%$ of the coupons are still missing.
\end{wsolution}
\end{exercisebox}

\begin{exercisebox}
\begin{pproblem}[The numbers for $n = 365$]\label{pp:bday365}
    \leavevmode
    \begin{enumerate}[label=(\alph*)]
        \item How many people make the expected number of collisions at least $1$?
        \item How many make the probability of a collision at least $\frac{1}{2}$, by the chain-rule
              bound?
        \item How many does the Chebyshev proof need to guarantee probability $\frac{1}{2}$?
    \end{enumerate}
\end{pproblem}
\tcblower
\begin{wsolution}
    \uses{\scr{thm:birthday-0903}}
    (a) We need $m(m-1) \geq 2n = 730$. Since $27 \cdot 26 = 702$ and $28 \cdot 27 = 756$, the
    answer is $m = 28$, with $\expectation{X} = \sfrac{756}{730} \approx 1.04$.

    (b) We need $e^{-m(m-1)/730} \leq \frac{1}{2}$, that is $m(m-1) \geq 730 \ln 2 \approx 505.997$.
    Since $22 \cdot 21 = 462$ and $23 \cdot 22 = 506$, the answer is $m = 23$. The exact product for
    $m = 23$ is $0.4927$, so a collision has probability $0.507$.

    (c) Chebyshev gives $\prob{X = 0} \leq \sfrac{1}{\expectation{X}}$, so we need
    $\expectation{X} \geq 2$, that is $m(m-1) \geq 1460$. Since $38 \cdot 37 = 1406$ and
    $39 \cdot 38 = 1482$, the answer is $m = 39$. The condition $m \geq 2\sqrt{n} + 1 \approx 39.2$
    of the exercise gives $m = 40$.

    The three answers differ ($28$, $23$, $39$) but all are about $\sqrt{n}$. Part (a) alone proves
    nothing about the probability: $\expectation{X} \geq 1$ does not imply that
    $\prob{X \geq 1}$ is large. That step needs more than the expectation: either the direct
    product of part (b), or the variance, which is what the exercise uses with Chebyshev.
\end{wsolution}
\end{exercisebox}


\newweek{3}{September 7--11, 2026}

\newlecture{Tuesday, September 8, 2026}{Distributed Computing, Randomized Leader Election in a Complete
    Network, Random Walks, Law of Large Numbers}
\input{tex/lecture_2026-09-08}

\newlecture{Thursday, September 10, 2026}{Law of Large Numbers, Random Walks on Graphs, Transition Matrix,
    Stationary Distribution and Convergence}
\input{tex/lecture_2026-09-10}

\practiceweek{3}

\begin{exercisebox}
\begin{pproblem}[$p = 3 \ln n / n$]\label{pp:leader3}
    Redo the analysis with candidate probability $p = \sfrac{3 \ln n}{n}$ and the same $k$.
\end{pproblem}
\tcblower
\begin{wsolution}
    \uses{\scr{thm:leader-atleast-0908}; \scr{lem:disjoint-0908}; \scr{thm:leader-atmost-0908}; \scr{thm:complexity-0908}}
    \leavevmode
    \begin{enumerate}
        \item At least one leader: the candidate with the smallest id always wins at all its
              referees, so there is no leader only if there is no candidate.
              $\prob{\text{no candidate}} = (1-p)^n \leq e^{-pn} = e^{-3 \ln n} = \sfrac{1}{n^3}$.
        \item Two candidates with a common referee cannot both be leaders. Two referee sets are
              disjoint with probability at most $e^{-k^2/n} \leq e^{-\ln n} = \sfrac{1}{n}$, as
              before.
        \item A fixed pair is bad (both are candidates, with no common referee) with probability at
              most $\sfrac{p^2}{n}$. The expected number of bad pairs is at most
              \[\binom{n}{2} \frac{p^2}{n} \leq \frac{n^2}{2} \cdot \frac{9 \ln^2 n}{n^2} \cdot \frac{1}{n} = \frac{9 \ln^2 n}{2n}.\]
        \item By Markov,
              $\prob{\text{more than one leader}} \leq \prob{\text{at least one bad pair}} \leq \frac{9 \ln^2 n}{2n}$.
        \item Messages: at most $2k$ per candidate and $3 \ln n$ candidates in expectation, so at
              most $6k \ln n = \bigO{\sqrt{n} \log^{3/2} n}$.
    \end{enumerate}
    More candidates make ``no leader'' less likely, and ``two leaders'' and the message count
    slightly larger. The asymptotics do not change.
\end{wsolution}
\end{exercisebox}

\begin{exercisebox}
\begin{pproblem}[A triangle with a tail]\label{pp:triangle}
    Vertices $1, 2, 3$ form a triangle and vertex $4$ is joined to vertex $3$.
    \begin{enumerate}[label=(\alph*)]
        \item Write $P$.
        \item Find $\pi$ and verify it.
        \item Does the walk converge?
        \item Starting at vertex $4$, compute $p_1, p_2, p_3$.
    \end{enumerate}
\end{pproblem}
\tcblower
\begin{wsolution}
    \uses{\scr{lem:evolution-0910}; \scr{thm:stationary-0910}; \scr{thm:convergence-0910}}
    (a) The degrees are $2, 2, 3, 1$, so
    \[P = \begin{pmatrix}
            0            & \tfrac{1}{2} & \tfrac{1}{2} & 0            \\
            \tfrac{1}{2} & 0            & \tfrac{1}{2} & 0            \\
            \tfrac{1}{3} & \tfrac{1}{3} & 0            & \tfrac{1}{3} \\
            0            & 0            & 1            & 0
        \end{pmatrix}.\]
    Every row sums to $1$.

    (b) There are $m = 4$ edges, so
    $\pi = \left(\frac{2}{8}, \frac{2}{8}, \frac{3}{8}, \frac{1}{8}\right)
        = \left(\frac{1}{4}, \frac{1}{4}, \frac{3}{8}, \frac{1}{8}\right)$.
    Verify $\pi P = \pi$ entry by entry:
    \begin{itemize}
        \item Vertex $1$: $\pi_2 \cdot \frac{1}{2} + \pi_3 \cdot \frac{1}{3} = \frac{1}{8} + \frac{1}{8} = \frac{1}{4}$.
        \item Vertex $2$: $\pi_1 \cdot \frac{1}{2} + \pi_3 \cdot \frac{1}{3} = \frac{1}{8} + \frac{1}{8} = \frac{1}{4}$.
        \item Vertex $3$: $\pi_1 \cdot \frac{1}{2} + \pi_2 \cdot \frac{1}{2} + \pi_4 \cdot 1
                  = \frac{1}{8} + \frac{1}{8} + \frac{1}{8} = \frac{3}{8}$.
        \item Vertex $4$: $\pi_3 \cdot \frac{1}{3} = \frac{1}{8}$.
    \end{itemize}

    (c) Yes. The graph is connected, and it contains a triangle (an odd cycle), so it is not
    bipartite.

    (d) $p_0 = (0, 0, 0, 1)$.
    \begin{itemize}
        \item $p_1 = (0, 0, 1, 0)$: from $4$ the walk must go to $3$.
        \item $p_2 = \left(\frac{1}{3}, \frac{1}{3}, 0, \frac{1}{3}\right)$: from $3$ each
              neighbour has probability $\frac{1}{3}$.
        \item $p_3 = \left(\frac{1}{3} \cdot \frac{1}{2},\ \frac{1}{3} \cdot \frac{1}{2},\
                  \frac{1}{3} \cdot \frac{1}{2} + \frac{1}{3} \cdot \frac{1}{2} + \frac{1}{3} \cdot 1,\ 0\right)
                  = \left(\frac{1}{6}, \frac{1}{6}, \frac{2}{3}, 0\right)$.
    \end{itemize}
    Each $p_t$ sums to $1$, as it must (a row stochastic matrix preserves probability vectors).
\end{wsolution}
\end{exercisebox}


\newweek{4}{September 14--18, 2026}

\newlecture{Tuesday, September 15, 2026}{Qubits and Dirac Notation, Unitary Matrices, Quantum Gates and
    Circuits, Tensor Products, Entanglement}
\input{tex/lecture_2026-09-15}

\newlecture{Thursday, September 17, 2026}{Linear Algebra Review: Column Space, Rank, Linear Independence,
    Orthogonality}
\input{tex/lecture_2026-09-17}

\practiceweek{4}

\begin{exercisebox}
\begin{pproblem}[Five states]\label{pp:five}
    Which of these states are entangled? Prove each answer.
    \begin{enumerate}[label=(\alph*)]
        \item $\frac{1}{2}\left(\ket{00} + \ket{01} + \ket{10} + \ket{11}\right)$
        \item $\frac{1}{2}\left(\ket{00} - \ket{01} + \ket{10} - \ket{11}\right)$
        \item $\frac{1}{2}\left(\ket{00} + \ket{01} + \ket{10} - \ket{11}\right)$
        \item $\frac{1}{\sqrt{2}}\left(\ket{01} - \ket{10}\right)$
        \item $\sqrt{\sfrac{1}{3}}\, \ket{00} + \sqrt{\sfrac{2}{3}}\, \ket{01}$
    \end{enumerate}
\end{pproblem}
\tcblower
\begin{wsolution}
    \uses{\scr{thm:bell-0915}}
    Compare $ad$ with $bc$:
    \begin{center}
        \begin{tabular}{cccl}
            \toprule
            State & $ad$              & $bc$              & Answer                                                                      \\ \midrule
            (a)   & $\sfrac{1}{4}$    & $\sfrac{1}{4}$    & Product: $\ket{+}\ket{+}$                                                   \\
            (b)   & $-\sfrac{1}{4}$   & $-\sfrac{1}{4}$   & Product: $\ket{+}\ket{-}$                                                   \\
            (c)   & $-\sfrac{1}{4}$   & $\sfrac{1}{4}$    & Entangled                                                                   \\
            (d)   & $0$               & $-\sfrac{1}{2}$   & Entangled                                                                   \\
            (e)   & $0$               & $0$               & Product: $\ket{0}\left(\sqrt{\sfrac{1}{3}}\,\ket{0} + \sqrt{\sfrac{2}{3}}\,\ket{1}\right)$ \\ \bottomrule
        \end{tabular}
    \end{center}

    \emph{Full write-up for (c).} Suppose the state equals
    $\left(\alpha_0\ket{0} + \alpha_1\ket{1}\right)\left(\beta_0\ket{0} + \beta_1\ket{1}\right)$.
    Then $\alpha_0\beta_0 = \frac{1}{2}$, $\alpha_0\beta_1 = \frac{1}{2}$,
    $\alpha_1\beta_0 = \frac{1}{2}$ and $\alpha_1\beta_1 = -\frac{1}{2}$. The product of the first
    and the last is $\alpha_0\alpha_1\beta_0\beta_1 = -\frac{1}{4}$. The product of the middle two
    is $\alpha_0\alpha_1\beta_0\beta_1 = +\frac{1}{4}$. This is a contradiction, so the state is
    entangled.
\end{wsolution}
\end{exercisebox}

\begin{exercisebox}
\begin{pproblem}[Is the matrix unitary?]\label{pp:unitary}
    Decide for each matrix whether it is unitary.
    \[U = \frac{1}{\sqrt{2}} \begin{pmatrix} 1 & i \\ i & 1 \end{pmatrix},
        \qquad
        V = \begin{pmatrix} 1 & 1 \\ 0 & 1 \end{pmatrix},
        \qquad
        W = \frac{1}{\sqrt{2}} \begin{pmatrix} 1 & 1 \\ 1 & 1 \end{pmatrix}.\]
\end{pproblem}
\tcblower
\begin{wsolution}
    \uses{\scr{thm:unitary-0915}}
    Compute $U^\dagger U$, where $\dagger$ means: transpose, and conjugate every entry.
    \[U^\dagger U = \frac{1}{2} \begin{pmatrix} 1 & -i \\ -i & 1 \end{pmatrix} \begin{pmatrix} 1 & i \\ i & 1 \end{pmatrix}
        = \frac{1}{2} \begin{pmatrix} 1 - i^2 & i - i \\ -i + i & -i^2 + 1 \end{pmatrix}
        = \begin{pmatrix} 1 & 0 \\ 0 & 1 \end{pmatrix}.\]
    So $U$ is unitary. Do not forget the conjugate: without it the diagonal would be $1 + i^2 = 0$.

    $V$ is not unitary: $V\ket{1} = \ket{0} + \ket{1}$ has norm $\sqrt{2}$, and a unitary matrix
    preserves norms (Sep 15). $W$ is not unitary: $W\ket{0} = W\ket{1} = \ket{+}$, so two orthogonal
    states are sent to the same state; also $W\left(\ket{0} - \ket{1}\right) = 0$.

    \emph{Quick test.} A matrix is unitary exactly when its columns are orthonormal, with the
    complex inner product.
\end{wsolution}
\end{exercisebox}


\newweek{5}{September 21--25, 2026}

\newlecture{Thursday, September 24, 2026}{Deutsch's Algorithm, Phase Flip, Walsh--Hadamard Transform,
    Deutsch--Jozsa Algorithm}
\input{tex/lecture_2026-09-24}

\practiceweek{5}

\begin{exercisebox}
\begin{pproblem}[Deutsch--Jozsa with $n = 2$]\label{pp:dj}
    Let $f(x_1 x_2) = x_1 \oplus x_2$. Run the Deutsch--Jozsa algorithm on $f$. What is measured,
    and what does the algorithm conclude?
\end{pproblem}
\tcblower
\begin{wsolution}
    \uses{\scr{thm:dj-0924}; \scr{lem:wh-0924}}
    \leavevmode
    \begin{enumerate}
        \item Start in $\ket{00}\ket{1}$. Hadamards on all three qubits give
              $\frac{1}{2}\left(\ket{00} + \ket{01} + \ket{10} + \ket{11}\right)\ket{-}$.
        \item The query multiplies each $\ket{x}$ by $(-1)^{f(x)}$. Here $f(00) = 0$, $f(01) = 1$,
              $f(10) = 1$, $f(11) = 0$, so the state is
              $\frac{1}{2}\left(\ket{00} - \ket{01} - \ket{10} + \ket{11}\right)\ket{-}$.
        \item Hadamards on the first two qubits. The amplitude of $\ket{y}$ is
              $\frac{1}{4} \sum_{x} (-1)^{f(x) + x \cdot y}$.
        \item For $y = 00$ this is $\frac{1}{4}(1 - 1 - 1 + 1) = 0$. For $y = 11$ we have
              $x \cdot y = x_1 \oplus x_2 = f(x)$, so every term is $+1$ and the amplitude is $1$.
        \item The measurement gives $11$ with probability $1$. The outcome is not $00$, so $f$ is
              balanced.
    \end{enumerate}
    One quantum query was used. A deterministic classical algorithm needs $2^{n-1} + 1 = 3$ queries.

    \emph{Check with a constant function.} For $f \equiv 1$ the state after the query is
    $-\frac{1}{2} \sum_{x} \ket{x} = -\ket{+}\ket{+}$, and the Hadamards turn it into $-\ket{00}$.
    The measurement gives $00$, so $f$ is constant. The global sign $-1$ does not affect the
    measurement.
\end{wsolution}
\end{exercisebox}

\begin{exercisebox}
\begin{pproblem}[Deutsch on $f(x) = 1 - x$]\label{wp:deutsch}
    Run Deutsch's algorithm on $f(0) = 1$, $f(1) = 0$, writing the state after every gate.
\end{pproblem}
\tcblower
\begin{wsolution}
    \uses{\scr{ex:hback-0924}; \scr{thm:deutsch-0924}}
    \leavevmode
    \begin{enumerate}
        \item Start in $\ket{0}\ket{1}$. After $H \otimes H$:
              $\ket{+}\ket{-} = \frac{1}{\sqrt{2}}\left(\ket{0} + \ket{1}\right)\ket{-}$.
        \item The query multiplies $\ket{x}$ by $(-1)^{f(x)}$:
              $\frac{1}{\sqrt{2}}\left(-\ket{0} + \ket{1}\right)\ket{-} = -\ket{-}\ket{-}$.
        \item $H$ on the first qubit: $H\ket{-} = \ket{1}$, so the state is $-\ket{1}\ket{-}$.
        \item The first qubit is measured as $1$ with probability $1$. Indeed
              $f(0) \oplus f(1) = 1$: $f$ is balanced.
    \end{enumerate}
    The sign $-1$ is global and does not change the measurement.
\end{wsolution}
\end{exercisebox}


\newweek{6}{September 28 -- October 2, 2026}

\newlecture{Tuesday, September 29, 2026}{Walsh--Hadamard Transform Recap, Bernstein--Vazirani Algorithm,
    Simon's Algorithm}
\input{tex/lecture_2026-09-29}

\newlecture{Thursday, October 1, 2026}{Classical Query Complexity of Simon's Problem, Yao's Min-Max
    Principle, Linear Independence in Simon's Algorithm}
\input{tex/lecture_2026-10-01}

\practiceweek{6}

\begin{exercisebox}
\begin{pproblem}[Bernstein--Vazirani with $n = 3$]\label{pp:bv}
    Let $f(x) = s \cdot x \bmod 2$ with the hidden string $s = 101$, so $f(x) = x_1 \oplus x_3$.
    \begin{enumerate}[label=(\alph*)]
        \item How does a classical algorithm find $s$?
        \item Run the quantum algorithm.
    \end{enumerate}
\end{pproblem}
\tcblower
\begin{wsolution}
    \uses{\scr{thm:bv-0929}; \scr{lem:signsum-0929}}
    (a) Query the unit vectors: $f(100) = 1$, $f(010) = 0$, $f(001) = 1$. The answers are the bits
    of $s$, so $n = 3$ queries. Fewer cannot work, since each answer is one bit and $s$ has $n$ bits.

    (b) After the first Hadamards and the query, the first register is
    \[\frac{1}{\sqrt{8}} \sum_{x} (-1)^{s \cdot x} \ket{x}
        = \frac{1}{\sqrt{8}} \left(\ket{000} - \ket{001} + \ket{010} - \ket{011} - \ket{100} + \ket{101} - \ket{110} + \ket{111}\right).\]
    After the second Hadamards the amplitude of $\ket{y}$ is
    \[\frac{1}{8} \sum_{x} (-1)^{s \cdot x + y \cdot x} = \frac{1}{8} \sum_{x} (-1)^{(s \oplus y) \cdot x}
        = \begin{cases} 1 & \text{if } y = s, \\ 0 & \text{if } y \neq s, \end{cases}\]
    by the sum-of-signs claim (Sep 29): the sum is $2^n$ if $s \oplus y = 000$ and $0$ otherwise.
    So the measurement gives $y = 101 = s$ with probability $1$, after one query.
\end{wsolution}
\end{exercisebox}

\begin{exercisebox}
\begin{pproblem}[Simon with $n = 3$]\label{pp:simon}
    The function $f \colon \set{0,1}^3 \to \set{0,1}^3$ is given by the table below. Run Simon's
    algorithm and find $s$.
    \begin{center}
        \begin{tabular}{c|cccccccc}
            $x$    & 000 & 001 & 010 & 011 & 100 & 101 & 110 & 111 \\ \hline
            $f(x)$ & 101 & 010 & 000 & 110 & 000 & 110 & 101 & 010
        \end{tabular}
    \end{center}
\end{pproblem}
\tcblower
\begin{wsolution}
    \uses{\scr{thm:simon-run-0929}; \scr{thm:simon-0929}; \scr{thm:indep-1001}}
    \leavevmode
    \begin{enumerate}
        \item After the Hadamards and the query, the state is
              $\frac{1}{\sqrt{8}} \sum_{x} \ket{x}\ket{f(x)}$.
        \item Measure the second register. Suppose the outcome is $010$. The inputs with
              $f(x) = 010$ are $001$ and $111$, so the first register is
              $\frac{1}{\sqrt{2}}\left(\ket{001} + \ket{111}\right)$.
        \item Apply $H^{\otimes 3}$. The amplitude of $\ket{y}$ is
              \[\frac{1}{\sqrt{2}} \cdot \frac{1}{\sqrt{8}} \left((-1)^{001 \cdot y} + (-1)^{111 \cdot y}\right)
                  = \frac{1}{4} (-1)^{001 \cdot y} \left(1 + (-1)^{110 \cdot y}\right).\]
        \item This is $\pm\frac{1}{2}$ if $110 \cdot y = 0$ and $0$ if $110 \cdot y = 1$. So $y$ is
              one of $000, 001, 110, 111$, each with probability $\sfrac{1}{4}$. The state is
              $\frac{1}{2}\left(\ket{000} - \ket{001} + \ket{110} - \ket{111}\right)$.
        \item Every outcome satisfies $y_1 \oplus y_2 = 0$. The same holds for any other outcome in
              step 2, because every pair $\set{x, x \oplus s}$ has the same difference.
        \item Run the circuit $n - 1 = 2$ times. Suppose the outcomes are $y = 111$ and $y = 001$.
              They give the equations $s_1 \oplus s_2 \oplus s_3 = 0$ and $s_3 = 0$.
        \item Solve: $s_3 = 0$ and $s_1 = s_2$. The solutions are $000$ and $110$, so the nonzero
              candidate is $s' = 110$.
        \item Check with two classical queries: $f(000) = 101 = f(110)$. So $f$ is $2$-to-$1$ and
              $s = 110$.
    \end{enumerate}

    \emph{What can go wrong.} If the two outcomes are dependent (for example $000$ and $111$, or
    $001$ twice), the system has more than two solutions and we run again. Two outcomes are
    independent with probability $\frac{3}{4} \cdot \frac{1}{2} = \frac{3}{8}$: the first must be
    nonzero ($3$ of $4$), and the second must be outside $\set{000, y_1}$ ($2$ of $4$). This matches
    the bound of at least $\sfrac{1}{4}$ from Oct 1.

    \emph{If $f$ were $1$-to-$1$.} Step 2 would leave a single $\ket{x}$, every $y$ would have
    probability $\sfrac{1}{8}$, and the check in step 8 would fail, so the algorithm would answer
    $s = 000$.
\end{wsolution}
\end{exercisebox}

\begin{exercisebox}
\begin{pproblem}[Independence for $n = 4$]\label{wp:indep4}
    Simon's algorithm with $n = 4$ collects $y_1, y_2, y_3$. Find the probability that they are
    linearly independent if $f$ is $2$-to-$1$, and if $f$ is $1$-to-$1$. How many attempts are
    expected in the $2$-to-$1$ case?
\end{pproblem}
\tcblower
\begin{wsolution}
    \uses{\scr{thm:indep-1001}}
    \emph{$2$-to-$1$.} The $y$'s are uniform in $V = \set{y \suchthat s \cdot y = 0}$, of dimension
    $d = 3$. Part 1 of the exercise gives
    \[\left(1 - \frac{1}{2}\right)\left(1 - \frac{1}{4}\right)\left(1 - \frac{1}{8}\right)
        = \frac{1}{2} \cdot \frac{3}{4} \cdot \frac{7}{8} = \frac{21}{64} \approx 0.33 \geq \frac{1}{4}.\]
    Each factor is the probability that the next vector is outside the span of the earlier ones:
    the span of $i - 1$ independent vectors has $2^{i-1}$ of the $2^3$ elements of $V$.

    \emph{$1$-to-$1$.} Now $V = \set{0,1}^4$, $d = 4$, and we need $d - 1 = 3$ independent
    vectors. Part 2 gives
    $\frac{3}{4} \cdot \frac{7}{8} \cdot \frac{15}{16} = \frac{315}{512} \approx 0.62 \geq \frac{1}{2}$.

    \emph{Attempts.} Each attempt succeeds with probability $\sfrac{21}{64}$, independently, so the
    number of attempts is geometric with mean $\sfrac{64}{21} \approx 3.05$ (\scr{thm:geometric-0827}).
\end{wsolution}
\end{exercisebox}


\newweek{7}{October 5--9, 2026}

\newlecture{Tuesday, October 6, 2026}{Grover's Search Algorithm, Phase Flip, Diffusion Operator,
    Rotation Analysis}
\input{tex/lecture_2026-10-06}

\newlecture{Thursday, October 8, 2026}{Grover's Algorithm as Two Reflections, Number of Iterations,
    Random Number of Iterations, Quantum Amplitude Amplification}
\input{tex/lecture_2026-10-08}
\practiceweek{7}

\begin{exercisebox}
\begin{pproblem}[$N = 8$, $a = 1$, and overshooting]\label{pp:grover8}
    Find the best $k$, and the success probability for $k = 1, 2, 3$.
\end{pproblem}
\tcblower
\begin{wsolution}
    \uses{\scr{cor:hk-1008}; \scr{thm:iterations-1008}}
    \leavevmode
    \begin{enumerate}
        \item $\sin\theta = \sfrac{1}{\sqrt{8}} \approx 0.3536$, so
              $\theta \approx 0.3614$ rad $\approx 20.7^\circ$.
        \item $k = \floor{\sfrac{\pi}{(4 \cdot 0.3614)}} = \floor{2.17} = 2$.
        \item Plug in:
    \end{enumerate}
    \begin{center}
        \begin{tabular}{ccc}
            \toprule
            $k$ & Angle $(2k+1)\theta$ & Success probability                \\ \midrule
            $0$ & $20.7^\circ$         & $\sfrac{1}{8} = 0.125$             \\
            $1$ & $62.1^\circ$         & $\sfrac{25}{32} \approx 0.781$     \\
            $2$ & $103.5^\circ$        & $\sfrac{121}{128} \approx 0.945$   \\
            $3$ & $144.9^\circ$        & $\sfrac{169}{512} \approx 0.330$   \\ \bottomrule
        \end{tabular}
    \end{center}
    The best is $k = 2$ with probability about $0.945$. A third iteration rotates past $90^\circ$
    and the probability drops to $0.33$: more iterations are not better.
\end{wsolution}
\end{exercisebox}

\begin{exercisebox}
\begin{pproblem}[Inversion about the mean for $N = 8$, $a = 1$]\label{wp:mean8}
    Track the eight amplitudes through one Grover iteration and compare with
    $\sin^2(3\theta)$.
\end{pproblem}
\tcblower
\begin{wsolution}
    \uses{\scr{lem:diffusion-1006}; \scr{cor:hk-1008}}
    \leavevmode
    \begin{enumerate}
        \item Start: every amplitude is $\sfrac{1}{\sqrt{8}}$.
        \item $Z_f$ negates the solution: seven amplitudes $\sfrac{1}{\sqrt{8}}$ and one
              $-\sfrac{1}{\sqrt{8}}$.
        \item The mean is $\frac{1}{8} \cdot \frac{7 - 1}{\sqrt{8}} = \frac{3}{4\sqrt{8}}$.
        \item Each amplitude $\alpha$ becomes $2 \cdot \text{mean} - \alpha$. Non-solutions:
              $\frac{6}{4\sqrt{8}} - \frac{1}{\sqrt{8}} = \frac{1}{2\sqrt{8}}$. Solution:
              $\frac{6}{4\sqrt{8}} + \frac{1}{\sqrt{8}} = \frac{5}{2\sqrt{8}}$.
        \item Success probability $\frac{25}{4 \cdot 8} = \frac{25}{32} \approx 0.781$. Check of the
              norm: $7 \cdot \frac{1}{32} + \frac{25}{32} = 1$.
    \end{enumerate}
    With $\sin\theta = \sfrac{1}{\sqrt{8}}$: $\sin 3\theta = 3\sin\theta - 4\sin^3\theta =
    \frac{3}{\sqrt{8}} - \frac{4}{8\sqrt{8}} = \frac{5}{2\sqrt{8}}$, the same amplitude.
\end{wsolution}
\end{exercisebox}

\begin{exercisebox}
\begin{pproblem}[$p = 1/100$]\label{pp:amp100}
    A quantum algorithm $\mathcal{A}$ outputs a good answer with probability $\sfrac{1}{100}$, and we
    can check an answer with $f$. How many runs of $\mathcal{A}$ are needed classically, and how
    many with amplitude amplification?
\end{pproblem}
\tcblower
\begin{wsolution}
    \uses{\scr{thm:ampamp-1008}}
    \leavevmode
    \begin{enumerate}
        \item Classically: run, check, repeat. The number of runs is geometric with mean
              $\sfrac{1}{p} = 100$.
        \item $\sin\theta = \sqrt{\sfrac{1}{100}} = 0.1$, so $\theta \approx 0.1002$.
        \item $k = \floor{\sfrac{\pi}{(4 \cdot 0.1002)}} = \floor{7.84} = 7$.
        \item The success probability is $\sin^2(15 \cdot 0.1002) = \sin^2(1.503) \approx 0.995$.
        \item Cost: $2 \cdot 7 + 1 = 15$ applications of $\mathcal{A}$ or its inverse, and $7$
              queries to $f$.
    \end{enumerate}
    So $15$ applications instead of $100$. In general the cost is at most
    $\frac{\pi}{2\sqrt{p}} + 1 = \bigO{\sfrac{1}{\sqrt{p}}}$ against $\sfrac{1}{p}$.
\end{wsolution}
\end{exercisebox}



% AI use declaration (required by the course if AI tools were used).
\clearpage
\section*{AI Use Declaration}

Tool: Claude Code (Anthropic).

I used AI to transcribe my handwritten notes from class and convert them to \LaTeX. I used AI as a
verifier and to explore: I worked out the steps and the ways to solve the problems manually, and then
used AI to verify them. The figures and examples come from the partial examples done in class, and
some are AI generated for practice.

% Uncomment if something is cited with \cite{...}.
% \bibliographystyle{abbrv}
% \bibliography{references}
\end{document}
